%% notes-tex/multimodal/043-datasets-primer/sections/5-proposal.tex
%% SOURCE: the owner's direction of 2026-10-19 (notes/experiments.043-expression-proteome-morphology.md,
%% "Design pass") and the 2026-10-20 review pass (same note, "Review pass"); numbers from
%% results/data_shapes_043.json, results/store_census.json, results/cross_dataset_correlation.json,
%% results/within_dataset_ceilings.json and results/distinct_alleles.json, through the tables
%% written by experiments/043-expression-proteome-morphology/scripts/primer_043_tables.py. The
%% case-study numbers cite their experiment's records: 031 (memory 031-inductive-fold-program,
%% jobs 5022 to 5081; operator audit jobs 4917 and 4918; closure plan header), 035 (rounds 2,
%% 9 to 11, slurm 3034, 3009) and 038 (slurm 3374, 3700, 3701), 010 and 030
%% (notes-tex/trigenic/010-additive-baselines; job 1605; job 3902; W&B xy3xnpau).

\section{The plan: every gene-level readout trained together\secstatus{todo}}
\label{sec:proposal}

Every gene-level readout of the store trains in one model, and the question it answers
is an engineering one: does the model predict the transcriptome of a deletion it has not
seen, and of a double deletion from its two measured singles? Table~\ref{tab:plan} gives
the sets. An earlier version of this section proposed two separate training batches,
deletions and then isolates with segregants; the owner superseded it on 2026-10-19, and the
review pass of 2026-10-20 replaced the separate split draws by one fold file and put a
linear ladder and a decision rule ahead of any training run.

%% SOURCE: experiments/043-expression-proteome-morphology/scripts/primer_043_tables.py
%% -> tables/t11-plan.tex (record counts from results/data_shapes_043.json and store_census.json)
\begin{table}[htbp]
\centering
\footnotesize
\caption[]{\textbf{The sets of the plan.} Records count the raw store (full pass of
\file{data_shapes_043.py}, slurm job 5902). The fold file is not yet written, so a split's
count is its expected fold-0 share under the stratified draw: a tenth of each deletion
dataset's genes, a fifth of the isolates by strain name, of each segregant panel and of the
promoter variants. Splits A, C, D and E are carved out of training records; split B is a
dataset of its own and never trains.}
\label{tab:plan}
%% GENERATED by experiments/043-expression-proteome-morphology/scripts/primer_043_tables.py
%% SOURCE: results/store_census.json, results/cross_dataset_correlation.json,
%% results/within_dataset_ceilings.json and the 019 ceiling files named in
%% tables/provenance.json. Do not edit by hand.

\begin{tabular}{>{\raggedright\arraybackslash}p{30mm}>{\raggedright\arraybackslash}p{52mm}>{\raggedleft\arraybackslash}p{16mm}>{\raggedright\arraybackslash}p{72mm}}
\toprule
Set & Datasets & Records or held-out units, n (expected) & Role \\
\midrule
Training, every gene-level readout & Kemmeren 2014, Hu 2007, Jackson 2020, Nadal-Ribelles 2025, Messner 2023, Zelezniak 2018, Caudal 2024, Muenzner 2024 protein, Muenzner 2024 turnover, Boocock 2025, N'Guessan 2025, Vaishnav 2022 & 168{,}785 & trained together; fold 0 of each key space is carved out of these records, fold 1 validates \\
Split A, deletion genes & every deletion dataset & 140 & gene fold 0: about a tenth of Kemmeren's 1{,}400 genes outside the 82 pinned Sameith singles; also Messner 2023 about 448, Nadal-Ribelles 2025 about 306, Hu 2007 about 26, Zelezniak 2018 about 10, Ohya 2005 about 468 \\
Split B, Sameith doubles & Sameith 2015 Dm & 72 & every double; scored against the per-gene mean, the stronger single, the additive null and the fitted combination of its singles \\
Split C, isolates & Caudal 2024, Muenzner 2024 protein and turnover & 196 & isolate fold 0 of 979 strains: about 159 Muenzner and 11 turnover isolates, out of every readout \\
Split D, segregants & Boocock 2025, N'Guessan 2025 & 16{,}493 & fold 0 within each cross: about 15{,}903 Boocock and 590 N'Guessan records \\
Split E, promoters & Vaishnav 2022 & 14{,}232 & fold 0 of 71{,}158 promoter variants; linear evaluations first \\
Probe & Ohya 2005 & 4{,}678 & ridge from the frozen latent to the CalMorph features, raw and growth-residualized; never trained \\
Dropped & Sameith 2015 Sm & 82 & the same arrays as Kemmeren (82 of 82) \\
\bottomrule
\end{tabular}

\end{table}

\notesec{Training}
Twelve datasets train together: every dataset of the store except Sameith 2015 Sm,
whose 82 profiles are Kemmeren's own arrays (82 of 82 identical), Sameith 2015 Dm, which
is split B, and Ohya 2005, which is the probe. Vaishnav 2022 stays in the plan as the
promoter signal: one sort-seq level per promoter insert, read at the record level through
the reporter token of Sec.~\ref{sec:input-token}, whose promoter block is the frozen FUDT
upstream encoder on the construct window; its linear evaluations run first and its joint
entry waits on the Addgene 127546 sequence. Expression and protein readouts come off the
perturbed graph: a gene's value is read from that gene's own perturbed node state, and the
CLS state and the perturbed-gene summary $z_S$ are the record-level sites
(Sec.~\ref{sec:output-sites}). The low-agreement datasets train too, and their measured
agreement is what the arm without them reads: Nadal-Ribelles agrees with itself across
batches at a median $r$ of 0.043 against 0.022 for different genotypes and with Kemmeren
at 0.002; two Jackson 2020 constructions of one genotype agree at 0.129 against 0.095; Hu
2007 released no replicate and agrees with Kemmeren at 0.107 per genotype, on an X score
that is not a log ratio; Zelezniak 2018 agrees with Messner at 0.049 per genotype over 715
proteins though its own ceiling is 0.916; Muenzner 2024 turnover agrees with abundance at
$-$0.006 per isolate.

\notesec{Split A, deletion genes}
One fold file holds every deletion gene's fold; fold 0 is held out of every deletion
dataset at once, so a held-out gene is never seen as a perturbation on any readout, and
the same genes test Kemmeren, Messner, Hu, Nadal-Ribelles and Zelezniak alone and jointly.
Kemmeren holds out about 140 of its 1{,}400 eligible genes (the 82 Sameith singles are
pinned to training). The score is the median per-gene Pearson across those strains; the
per-strain Pearson is reported centered, because two different Kemmeren mutants already
agree at a raw 0.284 (Sec.~\ref{sec:within}). The bound is Kemmeren's own-culture ceiling,
0.729 over 114 mutants with two cultures; the bars are the frozen-token ridge and kNN of
Sec.~\ref{sec:output-linear} on the same genes.

\notesec{Split B, the Sameith doubles}
All 72 Sameith 2015 double deletions are held out and scored, centered, against their
measured profiles and against four rungs that need no model: the per-gene mean, the
stronger single, the additive null of the two Kemmeren singles (the paper's own expected
value) and a fitted combination of the two singles. Sameith's Methods defer to Kemmeren's,
so the doubles were profiled by the same procedure; no double has a replicate, and the
Kemmeren ceiling of 0.729 is the bound for the same arrays.

\notesec{Split C, the isolates}
Fold 0 of the isolate folds, about 196 of the 979 Caudal and Muenzner strains with about
159 Muenzner and 11 turnover isolates among them, is held out of every readout of those
isolates. The bar is kinship ridge (GBLUP) on the same isolates; the reference is how far
another readout of the same isolate agrees, Caudal mRNA against Muenzner protein at a
median per-isolate $r$ of 0.208, against 0.031 for Kemmeren mRNA against Messner protein on
deletions. Splits D (segregants, within cross) and E (promoters) come from the same file.

\notesec{The morphology probe}
Ohya 2005 CalMorph enters no training. Its 4{,}678 genotypes are run through the frozen
model and a ridge reads the 116 moving features from $h$, $z_S$, $\bar H$ and their
concatenation, raw and with the Costanzo single-mutant fitness regressed out, since growth
matrices alone predict morphology nearly as well as the proteome (019 screen). A ridge on
the measured inputs already reaches a median $r$ of 0.348 from Kemmeren expression
(1{,}437 shared genotypes), 0.278 from Messner protein (4{,}279) and 0.235 from Hu (259)
over the 116 moving features, against row-permutation nulls near zero, and Ohya's own
ceiling is 0.611 over 122 wild-type replicates on all 281 (Table~\ref{tab:morph}).
Sec.~\ref{sec:output-protocol} gives the bars and the two reads.

\notesec{The genotypes the model has not yet read}
Isolates and segregants need a gene-level genotype the \texttt{Perturbation} processor
does not build: marking every isolate perturbation sets 5{,}575 of 6{,}607 nodes, and a
segregant sets none. Sec.~\ref{sec:input} gives the replacement: each isolate as its own
reference context built from its allele embeddings, with an interim arm that marks loss of
function and absence on S288C, and the definite rule for segregants (a gene inherits the
non-reference parent only inside a block called at posterior 0.9 or above), with the
parent's real allele token. The posterior cut removes little on the pooled panels (a
median of 0\% of markers for Boocock pooled, 1.7\% for N'Guessan) and a quarter to a third
on Boocock's de novo cells, whose single cells carry medians of 2{,}000 to 4{,}000 UMIs;
whether those cells train is decision P1. The measured genotype signal across segregant
panels is small: nearest-genotype deviation $r$ 0.113 against 0.070 for random pairs
(Sec.~\ref{sec:within}).

\subsection{What the three case studies teach\secstatus{todo}}
\label{sec:lessons}

Three earlier campaigns trained the same cell graph transformer on other labels, and each
found a bar, a scoring trap and a protocol that the first 043 design did not carry.

\notesec{Betaxanthin, experiment 031}
The GO recipe with a rank term reached a fold-average held-out Spearman of 0.229 to 0.230
over six seeds (fold sd 0.003 to 0.006), 0.264 on the pinned fold; the annotation-free cell
(composite sequence token plus metabolic block) 0.108 on the fold average and 0.125 on the
pin. The bars were ridge on the deleted gene's raw ProtT5 token, 0.082 on the pin, and
forests on ProtT5 (0.150) and on GO plus pFBA (0.204) on the fold average. What made it
honest: five inductive folds over deletion genes with the external pin as fold 0, every arm
paired by seed against a named control, the rule of a gain of at least 0.02 with 5 of 6
seeds up (9 of 12 at twelve), and a fixed budget, since epoch selection on validation was
null ($-0.009$). The operator audit (jobs 4917, 4918) found the single-deletion score to be
a learned affine read of the deleted gene's own token, reproduced by a ridge from that row:
the inductive gap is in the representation, not the head. Split A is a single-deletion task
and the same audit applies, so the first number to produce is the ridge from the deleted
gene's frozen composite token to its profile; the first 043 design ran no such bar and
proposed three initializations by two split seeds with no fold structure, where 031 found
the pin's sampling sd (0.041) equal to the seed sd.

\notesec{Environment, experiments 035 and 038}
Nested ridge on raw FCFP4 counts reached a median centered Spearman per held-out compound
of 0.303 over 123 compound-evaluations (035, slurm 3034) and 0.359 over 96 on the second
build (038, slurm 3374); the transformer trained to saturation tied ridge on the first
build (every entry point's compound-level interval through zero) and sat below it on the
second (0.312 and 0.297; slurm 3700, 3701). The protocol that fixed the verdict: fit on the
whole non-test pool as ridge does, a fixed 50 epochs keeping the last, three seeds, paired
against ridge on the same held-out units, and the interval resampled over the 41 compounds
rather than the 123 evaluations. Its centering lesson: each side is centered by its own
training mean, since centering the prediction by the measured mean rewarded an untrained
model at 0.356 after 60 steps (slurm 3009). The per-strain Pearson across genes has the
same trap here: Kemmeren's different-mutant pairs agree at 0.284 (job 5893), the shared
deletion response, which a model that emits the training-mean profile earns. Every
per-strain Pearson is therefore reported centered, and the per-gene mean is the first rung.

\notesec{Trigenic, experiments 010 and 030}
The trigenic Pearson reached $0.454 \pm 0.006$ on the random split and fell to about 0.13
for every model on held-out query pairs, with the fitness-head model recovering 0.41 on one
seed. On the random split the ladder read 0.390 for the query pair alone, 0.400 for the
additive per-gene ridge, 0.425 to 0.427 for an embedding MLP and 0.438 to 0.455 for the
transformer; on the 030 pool a zero-parameter lower-order null read 0.577 and 0.566 and
boosted trees over the lower-order terms 0.743 and 0.700, above the joint transformer's
0.562 and 0.391 (job 3902). The per-entry dataset token (a 15-way one-hot projected to 8
dimensions at the readout) converged at a validation interaction Pearson of 0.5655 on one
seed. Ordering generalized and magnitudes did not (held-out calibration slope 0.24, job
1605), and the slope is a property of the population scored (0.88 on the whole test split
against 1.59 inside the predicted top 100). Split B carries the additive null but needs the
whole ladder of rungs that see only a double's components, and because the binned head's
point prediction is an expectation that shrinks toward the mean, Spearman sits beside
Pearson and the calibration slope is named by population.

\notesec{The unified protocol}
Four steps, in order. (1) The linear ladder on CPU from the frozen embeddings, on the same
held-out units the transformer will use (Sec.~\ref{sec:output-linear}); nothing trains on a
card until it is on file. (2) The rule, committed as \file{conf/decision\_rule.yaml} before
the first GPU run: primary statistic per split, the control arm, a gain of at least 0.02,
at least 5 of 6 paired seeds up, and a 95\% interval resampled over held-out strains (genes
for per-gene statistics) that excludes zero. (3) Training with the budget and window
declared at launch; the validation fold monitors the escape and never selects the epoch;
the test fold is read once. (4) Paired comparison by initialization seed, each arm against
its named control, at the unit of generalization: held-out deletion gene, isolate, double,
segregant or promoter, with both intervals computed and the unit-level one quoted.

\notesec{Decisions}
Table~\ref{tab:decisions} collects every choice the design leaves to the owner, with
two options, the measured fact that bears on it and the reconciled proposal with its
named alternatives. The I rows come from the input side (Sec.~\ref{sec:input-decisions}),
the O rows from the output side (Sec.~\ref{sec:output-decisions}), the P rows from putting
the plan together, the S rows from the protocol the review pass added, D1 from the data per
model part (Sec.~\ref{sec:data-per-part}) and E1 from the isolate embedding plan
(Sec.~\ref{sec:input-embeddings}).

%% SOURCE: experiments/043-expression-proteome-morphology/scripts/primer_043_tables.py
%% -> tables/t13-decisions-rows.tex (DECISIONS, numbers filled by facts() from the result files)
\begin{landscape}
\scriptsize
\setlength{\tabcolsep}{3pt}
\begin{longtable}{>{\raggedright\arraybackslash}p{6mm}>{\raggedright\arraybackslash}p{22mm}>{\raggedright\arraybackslash}p{44mm}>{\raggedright\arraybackslash}p{30mm}>{\raggedright\arraybackslash}p{56mm}>{\raggedright\arraybackslash}p{62mm}>{\raggedright\arraybackslash}p{9mm}}
\caption[]{\textbf{Owner decisions, one checklist.} Each row is one choice the design
cannot make alone. Option 1 is the reconciled main option; the proposed column states the
decision and names its alternatives; ``open'' marks a choice with no proposal. The
measured fact is the one its section cites. The last column is the section that carries
the full argument.}\label{tab:decisions}\\
\toprule
Id & Decision & Option 1 & Option 2 & Measured fact & Proposed, with alternatives & Sec. \\
\midrule
\endfirsthead
\toprule
Id & Decision & Option 1 & Option 2 & Measured fact & Proposed, with alternatives & Sec. \\
\midrule
\endhead
%% GENERATED by experiments/043-expression-proteome-morphology/scripts/primer_043_tables.py
%% SOURCE: the owner-decision lists of sections/input-design.tex and
%% sections/output-design.tex, condensed in DECISIONS. Do not edit by hand.
I1 & Base gene token & learned identity table (the 019 model) & frozen sequence encoder (ProtT5, nucleotide-transformer windows, FUDT) & every encoder is already a per-gene row on main; the 019 campaign used none; no comparison on expression measured & open & \ref{sec:input-decisions} \\
I2 & Absence and deletion & one change type & two change types & 135 natural absences per Caudal isolate; the similarity is measurable after training as the cosine of the two embeddings & two & \ref{sec:input-decisions} \\
I3 & Isolate marks in the mask & every type, $M\approx$ 5,565 & loss of function and absence, $M\approx$ 518 & operator cost scales with the largest $M$ in a batch; within-isolate mRNA-protein agreement 0.208 does not separate the two & open & \ref{sec:input-decisions} \\
I4 & Segregant gene membership & whole gene inside a block & gene midpoint inside a block & blocks end at marker positions and the crossover lies in the gap; straddling genes per segregant not counted & whole gene (the prototype) & \ref{sec:input-decisions} \\
I5 & Chromosomal position & on for every genotype class & on only when segregants enter & no run bears on it; nearest-genotype deviation $r$ 0.113 against 0.070 is the only segregant signal & open & \ref{sec:input-position} \\
I6 & Strain tokens & extra keys of the operator & added to the CLS row & at $M={}$1 an extra key makes $M={}$2 and changes the softmax, which confounds a comparison with 019 & extra keys (the prototype's layout) & \ref{sec:input-decisions} \\
I7 & Operator normalization at large M & softmax, magnitude fixed & sigmoid gate, magnitude grows with M & operator A run at $M={}$1 to 3 only; sigmoid gate 0.200 against 0.210, one seed; nothing at $M$ of 10 or more & open & \ref{sec:input-decisions} \\
I8 & Accessory ORFs and the Vaishnav reporter & encode each as an extra token & report and skip & 84 accessory presences per Caudal isolate; one reporter (YFP) for 71,158 Vaishnav records & skip accessory ORFs; encode the reporter, since Vaishnav trains & \ref{sec:input-decisions} \\
O1 & Heads & joined by readout type, 6 heads with 15 per-dataset affines & one head per dataset, 15 heads & same-type agreement 0.002 to 0.107 (expression), 0.049 (protein); an affine cannot raise a correlation; no joint training scored & joined; split as the control arm & \ref{sec:output-decisions} \\
O2 & Experiment-type token & flat id & ontology path & the served ontology is flat, 34 leaves under phenotype & flat now, path when the ontology carries one & \ref{sec:output-decisions} \\
O3 & Where the tokens enter & at the head only & on the CLS row, or as an attended token & experiment 035: $-0.007$, $-0.025$, $+0.008$ against ridge for the three forms, every CI through zero & head only; attended token as an arm & \ref{sec:output-decisions} \\
O4 & Loss & MSE & 19-quantile pinball & 019 morphology: MSE on $z_S$ escaped 12 of 12 runs, pinball 8 of 10 & MSE; pinball as an arm & \ref{sec:output-decisions} \\
O5 & Dataset weighting & equal per dataset & by ceiling, or by record count & Boocock and Vaishnav hold 87\% of the records; Kemmeren holds 1,482 & equal, agreement-weighted for the low-agreement sets; record count as an arm & \ref{sec:output-decisions} \\
O6 & Nadal-Ribelles sentinels & $|\log_2\mathrm{FC}|$ > 20 set to missing & keep the stored value & 20.2\% of cells are sentinels & missing & \ref{sec:output-decisions} \\
O7 & Jackson target & deletion effect against the HO control of the same environment & absolute $\log_2(\mathrm{CPM}+1)$ with the environment token & construction pairs agree at 0.129 against 0.095 on the effect; one genotype across environments at 0.010 & the effect & \ref{sec:output-decisions} \\
O8 & Split A leakage rule & held-out gene out of every deletion dataset & out of the expression datasets only & the arm without the natural-genome datasets measures the isolate and segregant side either way & every deletion dataset & \ref{sec:output-decisions} \\
O9 & Split C size & a fifth of the isolates, 159 & another fraction & no Muenzner ceiling; Caudal against Muenzner 0.208 per isolate is the reference & a fifth & \ref{sec:output-decisions} \\
O10 & Budget and window & 019's 500 epochs read over 400 to 499 & a longer budget & 019's 1,000-epoch runs scored lower over 900 to 999 than over 400 to 499 in 7 of 8 & declare with the launch & \ref{sec:output-decisions} \\
P1 & Boocock de novo cells & train all 79,123 single cells under the posterior rule & train the 393 pooled segregants only & de novo cells: median 2,000 to 4,000 UMIs, 24 to 35\% of markers below posterior 0.9 & open; the every-readout direction includes them & \ref{sec:proposal} \\
P2 & Split A against split B & draw split A from Kemmeren genes outside the 82 Sameith singles & draw split A from all 1,482 & split B's additive null needs both singles of each double in training; the overlap of a random draw is not counted & draw outside the Sameith singles & \ref{sec:proposal} \\
\bottomrule

\end{longtable}
\end{landscape}

\notesec{Questions for the owner}
\begin{enumerate}
\item \textbf{Isolates as reference contexts or the interim arm.} The main design encodes
each isolate as its own reference context (979 cached contexts, the operator carrying only
engineered edits) and needs the per-context encoder; the interim arm marks loss of function
and absence on S288C with allele tokens and needs nothing new. Which runs first?
\item \textbf{Addgene 127546.} May its plasmid sequence be fetched into
\file{torchcell-constructs/} with a provenance record, so the Vaishnav construct gets the
1{,}003 bp window the gene promoter encoder expects?
\item \textbf{ProtT5 or ESM2 for the first isolate pass.} 775{,}856 distinct proteins:
about 2.7 hours of ProtT5 on two cards under the stated throughput assumptions, ESM2 about
a third of that.
\item \textbf{Segregant parents.} The five non-BY parents carry Peter ids that are all
members of the gene-keyed allele tarball; confirm that their real allele tokens replace the
parent id on marked genes.
\item \textbf{Headline joint arm.} With or without Hu, Jackson and Nadal-Ribelles (Kemmeren
agreement 0.107, 0.043 to 0.076, 0.002)? They can only lower the Kemmeren score if their
signal does not transfer; the rule names its control before the first read.
\item \textbf{Clade-disjoint isolate split.} A second isolate fold file, disjoint by Peter
2018 clade, beside the random one, once the clade is on the isolate descriptor?
\item \textbf{The quoted per-strain statistic.} Confirm that the centered per-strain
Pearson, with the raw value and the 0.284 floor beside it, is the form quoted.
\end{enumerate}
